% preamble.tex --- shared by all logbook documentation files
% Engine: pdflatex

\documentclass[11pt, a4paper]{report}

% --- Encoding and fonts ---
\usepackage[T1]{fontenc}
\usepackage[utf8]{inputenc}
\usepackage{lmodern}

% --- Page geometry ---
\usepackage[
    top=2.5cm, bottom=2.5cm,
    left=2.8cm, right=2.8cm
]{geometry}

% --- Colours ---
\usepackage[dvipsnames]{xcolor}
\definecolor{pionorange}{RGB}{255,152,0}
\definecolor{codeblue}{RGB}{0,80,160}
\definecolor{notebg}{RGB}{255,248,230}
\definecolor{warningbg}{RGB}{255,235,235}
\definecolor{codebg}{RGB}{245,245,245}
\definecolor{tabhead}{RGB}{230,230,230}

% --- Hyperlinks ---
\usepackage[
    colorlinks=true,
    linkcolor=codeblue,
    urlcolor=codeblue,
    citecolor=codeblue,
    pdftitle={Logbook Entry Generator Documentation},
    pdfauthor={PionCT Experiment}
]{hyperref}

% --- Tables ---
\usepackage{booktabs}
\usepackage{array}
\usepackage{tabularx}
\usepackage{longtable}

% --- Code listings ---
\usepackage{listings}
\lstset{
    basicstyle=\small\ttfamily,
    backgroundcolor=\color{codebg},
    frame=single,
    framerule=0.4pt,
    rulecolor=\color{gray!40},
    breaklines=true,
    breakatwhitespace=true,
    showstringspaces=false,
    tabsize=2,
    xleftmargin=6pt,
    xrightmargin=6pt,
    aboveskip=8pt,
    belowskip=6pt,
}
\lstdefinelanguage{JSON}{
    basicstyle=\small\ttfamily,
    morestring=[b]",
    stringstyle=\color{codeblue},
    literate=
        *{:}{{{\color{black}{:}}}}{1}
         {\{}{{{\color{black}{\{}}}}{1}
         {\}}{{{\color{black}{\}}}}}{1}
         {[}{{{\color{black}{[}}}}{1}
         {]}{{{\color{black}{]}}}}{1},
}
\lstdefinelanguage{bash}{
    basicstyle=\small\ttfamily,
    morecomment=[l]{\#},
    commentstyle=\color{gray},
}
\lstdefinelanguage{javascript}{
    basicstyle=\small\ttfamily,
    keywords={const,let,var,function,return,if,else,for,await,async,true,false},
    keywordstyle=\color{codeblue},
    morestring=[b]',
    morestring=[b]",
    stringstyle=\color{BrickRed},
    morecomment=[l]{//},
    commentstyle=\color{gray},
}

% --- Inline code ---
\newcommand{\code}[1]{\texttt{\small#1}}

% --- Note / Warning boxes ---
\usepackage{tcolorbox}
\tcbuselibrary{skins, breakable}

\newtcolorbox{notebox}{
    enhanced, breakable,
    colback=notebg,
    colframe=pionorange,
    leftrule=4pt, rightrule=0pt, toprule=0pt, bottomrule=0pt,
    arc=0pt, outer arc=0pt,
    left=8pt, right=8pt, top=4pt, bottom=4pt,
    fontupper=\small,
}
\newtcolorbox{warningbox}{
    enhanced, breakable,
    colback=warningbg,
    colframe=red!70!black,
    leftrule=4pt, rightrule=0pt, toprule=0pt, bottomrule=0pt,
    arc=0pt, outer arc=0pt,
    left=8pt, right=8pt, top=4pt, bottom=4pt,
    fontupper=\small,
}

% --- Section styling ---
\usepackage{titlesec}
\titleformat{\chapter}[block]
    {\normalfont\Large\bfseries\color{pionorange}}
    {\thechapter.}{0.8em}{}
    [\vspace{-0.3em}\color{pionorange}\rule{\linewidth}{1.5pt}]
\titlespacing{\chapter}{0pt}{0pt}{12pt}

\titleformat{\section}
    {\normalfont\large\bfseries\color{black}}
    {\thesection}{0.8em}{}
\titleformat{\subsection}
    {\normalfont\normalsize\bfseries}
    {\thesubsection}{0.8em}{}

% --- Header / footer ---
\usepackage{fancyhdr}
\pagestyle{fancy}
\fancyhf{}
\fancyhead[L]{\small\color{gray}\leftmark}
\fancyhead[R]{\small\color{gray}Logbook Entry Generator}
\fancyfoot[C]{\small\thepage}
\renewcommand{\headrulewidth}{0.4pt}
\renewcommand{\chaptermark}[1]{\markboth{#1}{}}

% --- Misc ---
\usepackage{enumitem}
\setlist[itemize]{itemsep=2pt, topsep=4pt}
\setlist[enumerate]{itemsep=2pt, topsep=4pt}
\usepackage{parskip}
\setlength{\parskip}{6pt}
\setlength{\parindent}{0pt}


\begin{document}

\title{
    \vspace{2cm}
    {\color{pionorange}\rule{\linewidth}{2pt}}\\[0.6em]
    {\Huge\bfseries Logbook Entry Generator}\\[0.3em]
    {\Large Developer Guide}\\[0.4em]
    {\color{pionorange}\rule{\linewidth}{2pt}}\\[1em]
    {\large PionCT Experiment --- Jefferson Lab}
    \vspace{1cm}
}
\author{}
\date{\today}
\maketitle
\thispagestyle{empty}

\tableofcontents
\newpage

% ============================================================
\chapter{Overview}
% ============================================================

This document covers the internal architecture of the build system and
generated HTML tool, intended for anyone modifying or extending the code.

The system has two layers:

\begin{enumerate}
    \item \textbf{\code{build.py}} --- a Python 3.6+ script that reads
          \code{config.json} and generates
          \code{output/logbook\_generator.html} and (optionally)
          \code{output/apps\_script.gs}.
    \item \textbf{\code{output/logbook\_generator.html}} --- a self-contained
          single-page application that runs entirely in the browser. No
          framework, no bundler, no server.
\end{enumerate}

The HTML tool is not hand-edited after generation. All customisation happens
in \code{config.json}; re-running \code{build.py} regenerates the outputs.

% ============================================================
\chapter{\texttt{build.py} Architecture}
% ============================================================

The script is structured as a pipeline of pure functions, each responsible
for generating one fragment of the final HTML or Apps Script.

\section{Entry Point}

\begin{lstlisting}
main()
  +-- load_config()          validate config.json
  +-- compute_schema_hash()  detect Sheets-breaking field changes
  +-- embed_logo()           base64-encode logo or fall back to path
  +-- build_html()           assemble the full HTML string
  |     +-- build_add_run_entry_fields()      <-- form inputs
  |     +-- build_generate_html_run_table()   <-- generateHTML() JS
  |     +-- build_extract_run_data()          <-- extractRunData() JS
  |     +-- build_save_run_entries()          <-- saveToLocalStorage() JS
  |     +-- build_load_run_entries()          <-- loadFromLocalStorage() JS
  |     +-- build_import_run_entries()        <-- importFromHTML() JS
  |     +-- build_quick_links_form()          <-- form HTML
  |     +-- build_quick_links_save_js()       <-- save JS fragment
  |     +-- build_quick_links_load_js()       <-- load JS fragment
  |     +-- build_quick_links_clear_js()      <-- import clear JS fragment
  |     +-- build_quick_links_generate_html() <-- generateHTML() JS fragment
  |     `-- build_import_links()             <-- importFromHTML() link matcher
  `-- build_apps_script()    assemble the Apps Script string
\end{lstlisting}

All \code{build\_*} functions take the parsed config (or a subset of it) and
return a string. They have no side effects --- they do not write files or
access the DOM.

\section{Field Partitioning}

A core concept used throughout the script is splitting \code{run\_fields}
into subsets:

\begin{lstlisting}[language=Python]
user_fields       = [f for f in fields if f["source"] == "user"]
auto_form_fields  = [f for f in fields if f["source"] == "auto"
                                       and f["key"] != "date"]
form_fields       = [f for f in fields if f["source"] in ("user","auto")
                                       and f["key"] != "date"]
html_table_fields = [f for f in fields if f["source"] != "sheet-only"]
\end{lstlisting}

\begin{center}
\begin{tabularx}{\linewidth}{@{}lX@{}}
\toprule
\textbf{Subset} & \textbf{Used for} \\
\midrule
\code{form\_fields} &
    Inputs rendered in the \code{addRunEntry()} template; save/load loops \\
\code{html\_table\_fields} &
    Column headers and data cells in \code{generateHTML()} and
    \code{importFromHTML()} \\
All fields &
    \code{apps\_script.gs} HEADERS array and row construction \\
\bottomrule
\end{tabularx}
\end{center}

\code{date} is excluded from \code{form\_fields} because it is auto-generated
at send/generate time and never shown as an input; it is included in
\code{html\_table\_fields} so it appears in the generated logbook HTML table.

\section{CSS Class to JS Variable Name}

\code{js\_key(css\_class)} converts a CSS class to a JS variable name:

\begin{lstlisting}
run-hms-p      ->  strip "run-"  ->  hms-p  ->  camelCase  ->  hmsP
run-comments   ->  strip "run-"  ->  comments               ->  comments
date           ->  no prefix     ->  date                   ->  date
\end{lstlisting}

This name is used consistently across \code{extractRunData},
\code{saveToLocalStorage}, \code{loadFromLocalStorage}, and the Apps Script
row array.

\section{Schema Hash}

\code{compute\_schema\_hash()} hashes the ordered list of
\code{(key, sheet\_column)} tuples for all fields. This hash is stored in
\code{.schema\_hash} after each build. On the next build, if the hash has
changed and Sheets is enabled, a warning is printed.

The hash is also embedded as a comment in the generated HTML and Apps Script
so the deployed version can be identified.

\begin{notebox}
If \code{.schema\_hash} is missing (deleted or first build), the comparison
is skipped silently and a new file is written. No data is lost.
\end{notebox}

\section{Python Version Compatibility}

The script requires \textbf{Python 3.6+}. It uses only the standard library:
\code{base64}, \code{hashlib}, \code{json}, \code{mimetypes}, \code{os},
\code{sys}, \code{pathlib}.

No \code{removeprefix} (3.9+), no walrus operator (3.8+), no match
statements (3.10+).

% ============================================================
\chapter{Generated HTML Tool Architecture}
% ============================================================

\section{State Management}

All mutable state lives in one place: the browser's \code{localStorage},
under the key \code{\{experiment\_name\_lowercase\}LogbookData}. The key is
derived at build time from \code{experiment\_name} in the config.

\begin{lstlisting}
localStorage["pionctLogbookData"] = {
  summaryText, runPlanLink, checklistLink, ...   (fixed fields)
  sheetsWebAppUrl, sheetsSharedSecret,           (if Sheets enabled)
  additionalLinks: [ {text, url}, ... ],
  shiftEntries:    [ {time, entry}, ... ],
  runEntries:      [ {runNumber, startTime, ..., sentToSheet}, ... ],
  timestamp: "ISO 8601 string"
}
\end{lstlisting}

\code{saveToLocalStorage()} is called after every user interaction that
changes data. \code{loadFromLocalStorage()} is called on
\code{DOMContentLoaded}.

\section{Run Entry DOM Lifecycle}

\code{addRunEntry()} creates a new \code{.entry-item} div and appends it to
\code{\#runEntries}. All data for a run lives in the DOM of its
\code{.entry-item} --- there is no separate JS data structure.
Save/load/export all operate by querying the DOM:

\begin{lstlisting}[language=javascript]
document.querySelectorAll('#runEntries .entry-item').forEach(item => {
    const value = item.querySelector('.run-hms-p').value;
    ...
});
\end{lstlisting}

This means the DOM is the single source of truth during a session. The
\code{sentToSheet} state is stored as a \code{data-sent} attribute on the
\code{.sheet-led} span, read back by \code{saveToLocalStorage()}.

\section{Google Sheets Export Flow}

\begin{lstlisting}
sendRunToSheet(button)
  `-- extractRunData(item)         builds the JSON payload
  `-- postRunToSheet(runData)
        `-- fetch(webAppUrl, POST) sends JSON body
        `-- Apps Script doPost()  appends row to sheet
        `-- returns {ok, rowAppended}
  `-- updates LED + status text
  `-- saveToLocalStorage()         persists the green LED state
\end{lstlisting}

The \code{Content-Type} is set to \code{text/plain;charset=utf-8}
intentionally. Apps Script Web Apps do not respond correctly to CORS preflight
requests, which \code{application/json} would trigger. Using \code{text/plain}
sends the request as a ``simple request'' (no preflight) while the Apps Script
reads the raw body via \code{e.postData.contents} regardless of declared
content type.

\section{Import from HTML}

\code{importFromHTML()} uses \code{DOMParser} to parse the
previously-generated logbook HTML, then identifies tables by their header
text:

\begin{itemize}
    \item A table with headers \code{Time} and \code{Logbook Entry}
          $\rightarrow$ shift log table
    \item A table with headers \code{Run Number} and \code{Start Time}
          $\rightarrow$ run summary table
\end{itemize}

Fields are restored by column index. If the field list has changed between
the session that generated the HTML and the importing tool, columns may
misalign silently. This is a known limitation (see
Chapter~\ref{ch:limitations}).

% ============================================================
\chapter{Adding a New Field}
% ============================================================

\begin{enumerate}
    \item Add an entry to \code{run\_fields} in \code{config.json}:
\begin{lstlisting}[language=JSON]
{
  "key":          "run-my-field",
  "label":        "My Field",
  "type":         "text",
  "source":       "user",
  "sheet_column": "My Field",
  "sheet_align":  "left"
}
\end{lstlisting}
    \item Run \code{python3 build.py}.
    \item The field appears automatically in the form, the HTML table,
          save/load, extract, and (if Sheets enabled) the Apps Script row.
    \item If Sheets is enabled and the schema hash changed, redeploy the
          Apps Script.
\end{enumerate}

No changes to \code{build.py} or the HTML template are needed for a standard
\code{text} or \code{textarea} field with \code{source} of \code{user},
\code{auto}, or \code{sheet-only}.

% ============================================================
\chapter{Adding a New Field \texttt{source} Type}
% ============================================================

Currently supported: \code{user}, \code{auto}, \code{sheet-only}.

To add a new type (e.g.\ \code{computed} with a formula, or \code{dropdown}):

\begin{enumerate}
    \item Add handling in \code{build\_add\_run\_entry\_fields()} --- the
          function that generates the form input HTML for each field.
    \item Update \code{VALID\_SOURCES} in \code{load\_config()} to accept
          the new value.
    \item Decide which partition sets the new source belongs to (form, table,
          both) and update the partition logic at the top of
          \code{build\_html()}.
    \item If the new type changes what gets sent to the sheet, update
          \code{build\_apps\_script()} accordingly.
\end{enumerate}

% ============================================================
\chapter{Adding a New Field \texttt{type}}
% ============================================================

Currently supported: \code{text} (single-line input), \code{textarea}
(multi-line).

To add \code{select} (dropdown):

\begin{enumerate}
    \item Add \code{"select"} to \code{VALID\_TYPES} in
          \code{load\_config()}.
    \item In \code{build\_add\_run\_entry\_fields()}, add a branch that emits
          a \code{<select>} element. You will also need a new config key
          (e.g.\ \code{"options": [...]}) to carry the dropdown values.
    \item The save/load/export machinery reads \code{.value} from the DOM
          element and does not need changes --- \code{<select>} elements
          expose \code{.value} the same as inputs.
\end{enumerate}

% ============================================================
\chapter{Modifying the CSS / Visual Theme}
% ============================================================

The entire CSS is a single \code{<style>} block in the generated HTML,
generated verbatim by \code{build\_html()} in \code{build.py}. The colour
constants are:

\begin{center}
\begin{tabular}{@{}ll@{}}
\toprule
\textbf{Role} & \textbf{Value} \\
\midrule
Accent / headings        & \code{\#ff9800} (orange) \\
Secondary accent         & \code{\#ffb74d} (light orange) \\
Container background     & \code{linear-gradient(135deg, \#6b7176, \#5a6066)} \\
Page background          & \code{\#e6e6e6} \\
Input background         & \code{\#464B50} \\
Body text (on page bg)   & \code{\#2C3539} \\
Label text (on container)& \code{\#ffffff} \\
\bottomrule
\end{tabular}
\end{center}

To change the theme, modify the CSS string inside \code{build\_html()} in
\code{build.py} and rebuild. Do not edit the generated HTML directly --- it
will be overwritten on the next build.

% ============================================================
\chapter{Deployment Checklist}
% ============================================================

When releasing a new version:

\begin{itemize}
    \item[$\square$] Edit \code{config.json} as needed
    \item[$\square$] Run \code{python3 build.py}
    \item[$\square$] Check for the schema hash change warning
    \item[$\square$] If hash changed and Sheets is enabled: paste new
          \code{apps\_script.gs}, create a new deployment, update the Web
          App URL in the tool
    \item[$\square$] Test the generated HTML in a browser (add a run entry,
          send to sheet, reload and verify session restores, test Import from
          Previous HTML)
    \item[$\square$] Commit \code{config.json}, \code{build.py}, and
          \code{docs/} to the repository
    \item[$\square$] Do \textbf{not} commit \code{output/} (generated files)
          or \code{.schema\_hash} unless your team specifically wants to track
          generated artefacts
\end{itemize}

% ============================================================
\chapter{Known Limitations and Technical Debt}
\label{ch:limitations}
% ============================================================

\begin{itemize}
    \item \textbf{\code{importFromHTML} link matching} uses keyword heuristics
          on link text (first word of each quick-link label). If two links
          share a first word, one will shadow the other. A future fix would
          be to embed machine-readable metadata in the generated HTML.

    \item \textbf{No field validation} --- all inputs are \code{type="text"}.
          Numeric fields (beam current, momentum, etc.) accept arbitrary
          strings. Downstream errors in the sheet or generated HTML are
          possible if non-numeric values are entered.

    \item \textbf{Single Apps Script deployment URL} is stored in
          localStorage. If a second coordinator updates the deployment and
          the URL changes, each browser running the tool needs the new URL
          updated manually.

    \item \textbf{\code{totalTime} computation} treats any invalid
          \code{HH:MM} string as ``leave unchanged'' rather than clearing the
          field. A manually entered value survives a start/stop edit only if
          both times become invalid simultaneously.
\end{itemize}

\end{document}
